% 夜半用户手册 · 第 10 章：终端与自动化入口（中文）
% T-DOC-001 逐章交付（CH10 任务未到前预写）。
% 事实来源：README.md、docs/README.md、docs/tool-availability.md、
% crates/yeban-tui 存在性、cli.rs 开关。工具数不手抄文档数字。

\chapter{终端与自动化入口}
\label{ch:entry-points}

\classquote{天下同归而殊涂，一致而百虑。}{《周易·系辞下》}{图形界面、终端与 MCP 是不同入口，通向的是同一个工程。}





同一份工程可以从三个入口操作：桌面 GUI、终端 TUI、命令行 CLI，此外还有给
AI agent 用的 MCP 工具面。它们共用同一套工程模型与后端，所以做的是同一件事，
只是交互方式不同。这一章各给一节，告诉你什么时候用哪个。

\section{TUI}

夜半提供一个终端用户界面（TUI）入口，源码在 \texttt{crates/yeban-tui}。它给没有
图形桌面、或习惯在终端里工作的场景使用，不必开窗口；具体可用的面板与键位以
该入口自身的帮助为准，本章不在此手抄一份可能过期的清单。

\begin{table}[htbp]
\centering
\caption{三种入口各适合什么}
\label{tab:ch10-entry-compare}
\begin{tabular}{ll}
\hline
入口 & 适合的场景 \\
\hline
桌面 GUI & 日常编曲、写音符、混音、看界面 \\
TUI & 终端环境、无图形桌面、远程终端 \\
CLI & 批处理、脚本、导出、CI \\
MCP & AI agent 通过工具调用编辑与查询 \\
\hline
\end{tabular}
\end{table}

\section{CLI}

\subsection{先看帮助}

完整的开关与退出码以 \texttt{yeban-app --help} 为准；文档只是导览，不替代帮助
输出。几个常用入口：

\begin{description}
  \item[\texttt{--open <路径>}] 直接打开指定工程（GUI 或批处理都可接）。
  \item[\texttt{--export-wav}] 导出整轨 WAV，可配 \texttt{--loudness-target} 与 \texttt{--true-peak-ceiling} 做母带化；目标不可达时响亮失败，不静默降级。
  \item[\texttt{--export-stems <目录>}] 用与 MCP 同一份后端导出分轨；可配选轨开关只导点名的轨。
  \item[\texttt{--apply-ops}] 把一份编辑操作文件当作一个批次提交。
  \item[\texttt{--export-song-overview}] 导出整曲概览（能量/频率/密度/情绪四层与异常摘要）。
  \item[\texttt{--print-shortcuts}] 打印快捷键策略表（与 GUI「帮助 → 快捷键速查」同源）。
  \item[\texttt{--headless}] 无显示器机器上运行。
\end{description}

\subsection{feature 门控要知道}

有些开关由编译 feature 门控。缺省构建里与 MCP 相关的 \texttt{in-process-mcp}
默认开，但「编译进来」不等于「正在监听」：运行时还需要 \texttt{--enable-mcp-http}
或环境变量 \texttt{YEBAN\_MCP\_HTTP=1} 才会监听。用 \texttt{--no-default-features}
构建时，调用被门控的开关会得到点名 feature 的用法错误（退出码 2），而不是静默
忽略。实验性导出（如 \texttt{--export-als}、\texttt{--export-logic}）默认关闭，
给了开关才编译。

\section{MCP}

\subsection{两种形态，一份注册表}

MCP 有两形态，共用同一份 \texttt{TOOLS} 注册表与同一个领域实现：

\begin{description}
  \item[形态 A（库形态）] 由 \texttt{yeban-app} 进程内托管，AI agent 在正在运行的 DAW 上工作；传输为环回 HTTP（\textbf{仅} \texttt{127.0.0.1}）或进程内直调。
  \item[形态 B（二进制形态）] 独立 \texttt{yeban-mcp} 二进制走 stdio 批处理，无 GUI 环境也能用。
\end{description}

网络监听只绑环回地址，不绑 \texttt{0.0.0.0}——这是安全红线。\texttt{yeban-mcp}
自身的 \texttt{mcp-http} 绝不进它的默认 feature，只能由宿主显式转发打开。

\screenshot{ch10-mcp-modes}{MCP 两形态示意：GUI 进程内挂载走环回 HTTP，独立二进制走 stdio}

\begin{figure}[htbp]
\centering
\begin{tikzpicture}[node distance=1.2cm, >=stealth,
    box/.style={draw, rounded corners, align=center, minimum width=2.4cm, minimum height=0.9cm}]
  \node[box] (a) {同一份 TOOLS 注册表};
  \node[box, below left=of a] (b) {形态 A：app 内挂载};
  \node[box, below right=of a] (c) {形态 B：独立二进制};
  \node[box, below=of b] (d) {环回 HTTP 127.0.0.1};
  \node[box, below=of c] (e) {stdio 批处理};
  \draw[->] (a)--(b); \draw[->] (a)--(c); \draw[->] (b)--(d); \draw[->] (c)--(e);
\end{tikzpicture}
\caption{MCP 两形态共用一份注册表；网络监听只绑环回地址}
\label{fig:ch10-mcp-forms}
\end{figure}

\subsection{工具数量别信文档里的手抄值}

工具数以 \texttt{crates/yeban-mcp/src/tools.rs} 的三个常量与
\texttt{schemas/mcp-tools.schema.json} 的 enum 现读为准：文档里的数字曾多次过期，
历史手册里写的值互相矛盾。本章刻意不写具体数字。

部分工具依赖宿主注入的 provider（例如外部 SVS singing 服务）；在形态 B 里没有
注入方，调用这类工具会得到 \texttt{providerUnavailable} 契约的回复，而不是假装成功。

\section{帮助主题映射}

\begin{table}[htbp]
\centering
\caption{本章小节 $\to$ 应用内帮助锚点}
\label{tab:ch10-help}
\begin{tabular}{ll}
\hline
本章小节 & 帮助主题 \\
\hline
TUI & \texttt{entry-tui} \\
CLI：先看帮助 / 常用开关 & \texttt{entry-cli-switches} \\
CLI：feature 门控 & \texttt{entry-cli-features} \\
MCP：两种形态 & \texttt{entry-mcp-forms} \\
MCP：工具数量与 provider & \texttt{entry-mcp-tools} \\
\hline
\end{tabular}
\end{table}

\paragraph{本章学完能做什么} 按场景选对入口，知道 CLI 的常用开关与 feature 门控，并说清 MCP 两形态共用注册表、网络只绑环回地址。

\paragraph{动手练习} 用 \texttt{yeban-app --help} 找出导出分轨的确切开关，对同一工程分别用 GUI 与 CLI 导出一次分轨并对照回执；运行 \texttt{yeban-app --print-shortcuts} 与 GUI 速查面板对照；最后确认不带 \texttt{--enable-mcp-http} 时 MCP HTTP 端口没有在监听。
